\section{Formal Setup}\label{sec:formal-setup}

This section fixes the abstract, non-physical vocabulary used throughout the paper. We work with finite decision problems whose state spaces factor into coordinates, and we isolate the minimal terminology needed to state the later complexity results.

\subsection{Decision Problems with Coordinate Structure}

\begin{definition}[Decision Problem]\label{def:decision-problem}
A \emph{decision problem with coordinate structure} is a tuple $\mathcal{D} = (A, X_1, \ldots, X_n, U)$ where:
\begin{itemize}
\item $A$ is a finite set of \emph{actions} (alternatives)
\item $X_1, \ldots, X_n$ are finite \emph{coordinate spaces}
\item $S = X_1 \times \cdots \times X_n$ is the \emph{state space}
\item $U : A \times S \to \mathbb{Q}$ is the \emph{utility function}
\end{itemize}
\end{definition}

\begin{definition}[Projection]\label{def:projection}
For state $s = (s_1, \ldots, s_n) \in S$ and coordinate set $I \subseteq \{1, \ldots, n\}$:
\[
s_I := (s_i)_{i \in I}
\]
is the \emph{projection} of $s$ onto coordinates in $I$.
\end{definition}

\begin{definition}[Optimizer Map]\label{def:optimizer}
For state $s \in S$, the \emph{optimal action set} is:
\[
\Opt(s) := \arg\max_{a \in A} U(a, s) = \{a \in A : U(a,s) = \max_{a' \in A} U(a', s)\}
\]
\end{definition}

\subsection{Sufficiency and Relevance}

\begin{definition}[Sufficient Coordinate Set]\label{def:sufficient}
A coordinate set $I \subseteq \{1, \ldots, n\}$ is \emph{sufficient} for decision problem $\mathcal{D}$ if:
\[
\forall s, s' \in S: \quad s_I = s'_I \implies \Opt(s) = \Opt(s')
\]
Equivalently, the optimal action depends only on coordinates in $I$.
\end{definition}

\begin{proposition}[Empty-Set Sufficiency Equals Constant Decision Boundary]\label{prop:empty-sufficient-constant}
The empty set $\emptyset$ is sufficient if and only if $\Opt(s)$ is constant over the entire state space.
\end{proposition}

\begin{proof}
If $\emptyset$ is sufficient, then every two states agree on their empty projection, so sufficiency forces $\Opt(s)=\Opt(s')$ for all $s,s' \in S$. The converse is immediate.
\end{proof}

\begin{proposition}[Insufficiency Equals Counterexample Witness]\label{prop:insufficiency-counterexample}
Coordinate set $I$ is not sufficient if and only if there exist states $s,s' \in S$ such that $s_I=s'_I$ but $\Opt(s)\neq\Opt(s')$.
\end{proposition}

\begin{proof}
This is the direct negation of Definition~\ref{def:sufficient}.
\end{proof}

\begin{definition}[Minimal Sufficient Set]\label{def:minimal-sufficient}
A sufficient set $I$ is \emph{minimal} if no proper subset $I' \subsetneq I$ is sufficient.
\end{definition}

\begin{definition}[Relevant Coordinate]\label{def:relevant}
Coordinate $i$ is \emph{relevant} if there exist states that differ only at coordinate $i$ and induce different optimal-action sets:
\[
i \text{ is relevant}
\iff
\exists s,s' \in S:\;
\Big(\forall j \neq i,\; s_j=s'_j\Big)
\ \wedge\
\Opt(s)\neq\Opt(s').
\]
\end{definition}

\begin{proposition}[Minimal-Set/Relevance Equivalence]\label{prop:minimal-relevant-equiv}
For any minimal sufficient set $I$ and any coordinate $i$,
\[
i \in I \iff i \text{ is relevant}.
\]
Hence every minimal sufficient set is exactly the relevant-coordinate set.
\end{proposition}

\begin{proof}
If $i \in I$ were not relevant, then changing coordinate $i$ while keeping all others fixed would never alter the optimal-action set, so removing $i$ would preserve sufficiency, contradicting minimality. Conversely, every sufficient set must contain each relevant coordinate, because omitting a relevant coordinate leaves a witness pair with identical $I$-projection but different optimal-action sets.
\end{proof}

\begin{definition}[Exact Relevance Identifiability]\label{def:exact-identifiability}
For a decision problem $\mathcal{D}$ and candidate set $I$, we say that $I$ is \emph{exactly relevance-identifying} if
\[
\forall i \in \{1,\ldots,n\}:\quad i \in I \iff i \text{ is relevant for } \mathcal{D}.
\]
\end{definition}

\begin{definition}[Structural Rank]\label{def:srank}
The \emph{structural rank} of a decision problem is the number of relevant coordinates:
\[
\mathrm{srank}(\mathcal{D}) := \left|\{i \in \{1,\ldots,n\}: i \text{ is relevant}\}\right|.
\]
\end{definition}

\subsection{Optimizer Quotient}

\begin{definition}[Decision Equivalence]\label{def:decision-equiv}
States $s,s' \in S$ are \emph{decision-equivalent} if they induce the same optimal-action set:
\[
s \sim_{\Opt} s' \iff \Opt(s)=\Opt(s').
\]
\end{definition}

\begin{definition}[Decision Quotient]\label{def:decision-quotient}
The \emph{decision quotient} (equivalently, optimizer quotient) of $\mathcal{D}$ is the quotient object
\[
Q_{\mathcal{D}} := S / {\sim_{\Opt}}.
\]
Its equivalence classes are exactly the maximal subsets of states on which the optimal-action set is constant.
\end{definition}

\begin{proposition}[Optimizer Quotient as Coimage/Image Factorization]\label{prop:optimizer-coimage}
The decision quotient $Q_{\mathcal{D}}$ is canonically in bijection with $\operatorname{range}(\Opt)$. Equivalently, in \textbf{Set}, it is the coimage/image factorization of the optimizer map.
\end{proposition}

\begin{proof}
Two states are identified by $\sim_{\Opt}$ exactly when they have the same optimizer value. Hence quotient classes are in one-to-one correspondence with attained values of $\Opt$.
\end{proof}

\begin{proposition}[Sufficiency Characterization]\label{prop:sufficiency-char}
Coordinate set $I$ is sufficient if and only if the optimizer map factors through the projection $\pi_I : S \to S_I$. Equivalently, there exists a function $\overline{\Opt}_I : S_I \to \mathcal{P}(A)$ such that
\[
\Opt = \overline{\Opt}_I \circ \pi_I.
\]
\end{proposition}

\begin{proof}
If $I$ is sufficient, define $\overline{\Opt}_I(x)$ to be $\Opt(s)$ for any state $s$ with $s_I = x$; this is well defined exactly because sufficiency makes $\Opt$ constant on each projection fiber. Conversely, any such factorization forces states with equal $I$-projection to have equal optimal-action sets.
\end{proof}

\subsection{Computational Model and Input Encoding}\label{sec:encoding}

The same decision question yields different computational problems under different input models. We therefore fix the encoding regime before stating any complexity classification.

\paragraph{Succinct encoding (primary for hardness).}
An instance is encoded by:
\begin{itemize}
\item a finite action set $A$ given explicitly,
\item coordinate domains $X_1,\ldots,X_n$ given by their sizes in binary,
\item a Boolean or arithmetic circuit $C_U$ that on input $(a,s)$ outputs $U(a,s)$.
\end{itemize}
The input length is
\[
L = |A| + \sum_i \log |X_i| + |C_U|.
\]
Unless stated otherwise, all class-theoretic hardness results in this paper are measured in $L$.

\paragraph{Explicit-state encoding.}
The utility function is given as a full table over $A \times S$. The input length is
\[
L_{\mathrm{exp}} = \Theta(|A||S|)
\]
up to the bitlength of the utility entries. Polynomial-time claims stated in terms of $|S|$ use this explicit-state regime.

\paragraph{Query-access regime.}
The solver is given oracle access to $\Opt(s)$ or to the utility values needed to reconstruct it. Complexity is then measured by query count, optionally paired with per-query evaluation cost. This regime separates access obstruction from full-table availability and from succinct circuit input length.

\begin{remark}[Decision Questions vs Decision Problems]\label{rem:question-vs-problem}
The predicate ``is coordinate set $I$ sufficient for $\mathcal{D}$?'' is an encoding-independent mathematical question. A \emph{decision problem} arises only after fixing how $\mathcal{D}$ and $I$ are represented. Complexity classes are properties of these encoded problems, not of the abstract predicate in isolation.
\end{remark}
